% Домашнее задание № 1 по математическому анализу для факультета МТ
% Лаборатория математики ФН1
% Компиляция: pdflatex dz-01-mt-usloviya.tex

\documentclass[12pt, a4paper]{article}

\usepackage[T2A]{fontenc}
\usepackage[utf8]{inputenc}
\usepackage[russian]{babel}
\usepackage[left=25mm, right=20mm, top=20mm, bottom=20mm]{geometry}
\usepackage{amsmath, amssymb}
\usepackage{enumitem}

\pagestyle{plain}
\setlength{\parindent}{1.25cm}

\begin{document}

\begin{center}
  {\large\bfseries Домашнее задание № 1}

  \vspace{0.3em}
  {\large Производная в прикладных задачах}

  \vspace{0.5em}
  {\small Математический анализ · факультет «Машиностроительные технологии»\\
  1 курс, 1 семестр · Лаборатория математики ФН1}
\end{center}

\vspace{1em}

\noindent
Задание составлено для программы факультета МТ: меньше теоретических
доказательств, больше приложений к задачам механики и технологии.
Номер варианта совпадает с номером студента в журнале группы.

\section*{Задача 1. Скорость и ускорение}

Точка движется по закону $s(t) = a t^3 - b t^2 + c t$ (путь в метрах,
время в секундах). Найти скорость и ускорение в момент $t = 2$~с,
определить моменты времени, когда точка останавливается.

\section*{Задача 2. Наибольшая производительность}

Производительность станка описывается функцией
\[
P(v) = \frac{k v}{v^2 + m},
\]
где $v$~--- скорость подачи. Найти скорость, при которой производительность
максимальна, и убедиться, что это именно максимум.

\section*{Задача 3. Оптимальные размеры}

Требуется изготовить закрытый цилиндрический бак объёма $V$ с наименьшим
расходом материала. Найти отношение высоты к диаметру и объяснить,
почему на практике оно часто нарушается.

\section*{Задача 4. Приближённые вычисления}

С помощью дифференциала вычислить приближённо $\sqrt[3]{p}$ и оценить
погрешность. Сравнить с точным значением, полученным на калькуляторе.

\section*{Задача 5. Скорость изменения}

Стороны прямоугольного бака меняются со скоростями $\dot a$ и $\dot b$.
Найти скорость изменения площади основания в момент, когда $a = a_0$,
$b = b_0$.

\section*{Задача 6. Исследование функции}

Исследовать функцию $f(x) = x e^{-x^2}$ и построить график.
Указать точки экстремума и перегиба.

\vspace{1.2em}

\section*{Требования к оформлению}

\begin{itemize}[nosep]
  \item в прикладных задачах обязательно указывать единицы измерения;
  \item ответ выделяется отдельной строкой;
  \item в задачах 2 и 3 требуется проверка характера экстремума;
  \item график строится на координатной сетке.
\end{itemize}

\vspace{0.7em}
\noindent\textbf{Сроки.} Выдача~--- 4-я учебная неделя,
сдача~--- \textbf{до конца 10-й недели}.

\end{document}
